% ============================================================
%  Chapter 6 --- Query Optimiser
% ============================================================
\chapter{Query Optimiser}
\label{ch:query-optimizer}

\section{The Role of a Query Optimiser}

A query optimiser transforms a \emph{logical query plan} (the AST from the parser) into an efficient \emph{physical execution plan}.  The optimiser's job is to find the plan with the lowest estimated cost among the many equivalent plans that produce the same result~\cite{selinger1979access}.

\begin{conceptbox}[Why Optimisation Matters]
The same SQL query can be executed in vastly different ways.  For a three-table join with a filter, the number of possible execution orders is $3! = 6$, and each order can use different join algorithms.  A query with 10 tables has $10! \approx 3.6$ million orderings.  The optimiser must prune this space efficiently.
\end{conceptbox}

\section{Query Processing in \shiftdb{}}

\shiftdb{} uses a \emph{direct execution} model: the executor walks the AST and immediately produces results.  There is no separate optimiser pass, but the \code{EXPLAIN} command generates a logical query plan for inspection.

\begin{figure}[H]
\centering
\begin{tikzpicture}[node distance=0.7cm and 0.6cm]
  \node[component, minimum width=2.4cm] (ast) {AST};
  \node[component, minimum width=2.4cm, right=of ast, fill=gray!10, draw=gray, dashed] (opt) {Optimiser\\(future)};
  \node[component, minimum width=2.4cm, right=of opt] (exec) {Executor};
  \node[component, minimum width=2.4cm, right=of exec] (result) {Query\\Result};

  \draw[arr] (ast) -- (opt);
  \draw[arr] (opt) -- (exec);
  \draw[arr] (exec) -- (result);
\end{tikzpicture}
\caption{Current query execution path.  The optimiser stage is a logical placeholder.}
\label{fig:optimizer-path}
\end{figure}

\section{EXPLAIN Output}

The \code{EXPLAIN} command generates an operator tree for a SELECT statement:

\begin{lstlisting}[language=SQL, caption={Example EXPLAIN output.}]
EXPLAIN SELECT c.name, o.total_price
FROM customers c
INNER JOIN orders o ON c.id = o.customer_id
WHERE c.city = 'NYC'
ORDER BY o.total_price DESC
LIMIT 10;
\end{lstlisting}

\begin{table}[H]
\centering
\caption{Sample EXPLAIN result.}
\begin{tabular}{c l l l}
\toprule
\textbf{id} & \textbf{operation} & \textbf{table} & \textbf{details} \\
\midrule
1 & TABLE SCAN   & customers & Full table scan \\
2 & INNER JOIN   & orders    & Nested loop join \\
3 & FILTER       &           & WHERE clause evaluation \\
4 & SORT         &           & In-memory sort \\
5 & LIMIT        &           & Limit 10 \\
\bottomrule
\end{tabular}
\end{table}

\section{Classical Optimiser Architecture}

Although \shiftdb{} uses direct execution, understanding the classical optimiser is essential for interviews.

\subsection{Logical Plan $\to$ Physical Plan}

\begin{figure}[H]
\centering
\begin{tikzpicture}[node distance=1cm and 0.8cm]
  \node[component, minimum width=3cm] (lp) {Logical Plan\\(relational algebra)};
  \node[component, minimum width=3cm, right=2cm of lp] (rew) {Rewrite\\Rules};
  \node[component, minimum width=3cm, right=2cm of rew] (pp) {Physical Plan\\(algorithms)};
  \node[component, minimum width=3cm, below=of rew] (cost) {Cost\\Estimator};

  \draw[arr] (lp) -- (rew);
  \draw[arr] (rew) -- (pp);
  \draw[arr] (cost) -- (rew);
\end{tikzpicture}
\caption{Classical query optimisation pipeline.}
\end{figure}

\subsection{Key Rewrite Rules}

\begin{table}[H]
\centering
\caption{Common logical optimisation rules.}
\begin{tabular}{l p{12cm}}
\toprule
\textbf{Rule} & \textbf{Description} \\
\midrule
Predicate Pushdown    & Move filters (\code{WHERE}) as close to the data source as possible \\
Projection Pushdown   & Only read columns that are actually needed \\
Join Reordering       & Choose the join order with the smallest intermediate result sets \\
Constant Folding      & Evaluate constant expressions at compile time \\
Subquery Unnesting    & Convert correlated subqueries to joins \\
\bottomrule
\end{tabular}
\end{table}

\subsection{Cost Estimation}

The cost model estimates the expense of each plan using:
\begin{itemize}[nosep]
  \item \textbf{Cardinality estimates} --- expected number of output rows.
  \item \textbf{Selectivity} --- fraction of rows surviving a predicate.
  \item \textbf{I/O cost} --- number of page reads/writes.
  \item \textbf{CPU cost} --- number of comparisons, hashes, etc.
\end{itemize}

\begin{conceptbox}[Selectivity Estimation]
For a predicate \code{col = value}:
\[
  \text{selectivity} = \frac{1}{|\text{distinct values of col}|}
\]
For range predicates, databases use histograms (equi-width or equi-depth) to estimate the fraction of qualifying rows.
\end{conceptbox}

\section{Join Algorithms}

\shiftdb{} implements a \textbf{nested loop join} for all join types.  Production databases offer additional algorithms:

\begin{table}[H]
\centering
\caption{Join algorithm comparison.}
\begin{tabular}{l c p{8cm}}
\toprule
\textbf{Algorithm} & \textbf{Complexity} & \textbf{When Used} \\
\midrule
Nested Loop        & $\bigO{n \cdot m}$ & Small tables or indexed inner table \\
Block Nested Loop  & $\bigO{n \cdot m / B}$ & When buffer can hold blocks \\
Sort-Merge Join    & $\bigO{n\log n + m\log m}$ & Both sides sorted or need ORDER BY \\
Hash Join          & $\bigO{n + m}$ & Equality joins, sufficient memory \\
\bottomrule
\end{tabular}
\end{table}

\subsection{Nested Loop Join in \shiftdb{}}

\begin{lstlisting}[style=cpp, caption={Simplified nested loop join logic.}]
for (auto& leftRow : rows) {
    for (auto& rightRow : joinRows) {
        Row combined = leftRow;
        combined.insert(combined.end(),
                       rightRow.begin(), rightRow.end());
        // Build evaluation context and test ON condition
        if (isTruthy(evaluateExpr(join.onCondition, ctx, {})))
            joinedRows.push_back(combined);
    }
}
\end{lstlisting}

\section{Query Execution Strategies}

\begin{table}[H]
\centering
\caption{Execution models.}
\begin{tabular}{l p{9cm} l}
\toprule
\textbf{Model} & \textbf{How It Works} & \textbf{Used By} \\
\midrule
Volcano/Iterator  & Each operator has \code{next()}: pull one tuple at a time & PostgreSQL \\
Materialisation   & Each operator produces its full result before passing up & \shiftdb{} \\
Vectorised        & Process a batch of tuples per \code{next()} call & DuckDB \\
Push-based        & Data is pushed from source toward consumer & HyPer, Umbra \\
\bottomrule
\end{tabular}
\end{table}

\shiftdb{} uses \textbf{materialisation}: \code{scanAll()} loads all rows, the WHERE filter creates a new vector, and GROUP BY collects into a \code{std::map}.  This is simple but memory-intensive for large tables.

\section{Execution Steps in \code{executeSelect}}

\begin{enumerate}[nosep]
  \item \textbf{Base scan}: \code{table->scanAll()} retrieves all rows.
  \item \textbf{JOIN}: nested loop join appends columns from joined tables.
  \item \textbf{WHERE}: evaluates predicate per row, produces filtered vector.
  \item \textbf{GROUP BY / Aggregates}: groups rows by key, computes \code{COUNT}/\code{SUM}/\code{AVG}/\code{MIN}/\code{MAX}.
  \item \textbf{HAVING}: filters groups after aggregation.
  \item \textbf{ORDER BY}: \code{std::sort} with custom comparator.
  \item \textbf{OFFSET / LIMIT}: slices the result vector.
  \item \textbf{Projection}: evaluates SELECT expressions, builds output columns.
  \item \textbf{DISTINCT}: removes duplicate rows.
\end{enumerate}

\begin{interviewbox}[Interview Corner]
\textbf{Q:} What is the difference between a \emph{logical} and \emph{physical} query plan?\\[4pt]
\textbf{A:} A logical plan specifies \emph{what} operations to perform (filter, join, project) using relational algebra.  A physical plan specifies \emph{how} --- which algorithm to use (hash join vs.\ sort-merge) and which indexes to leverage.

\vspace{6pt}
\textbf{Q:} Explain predicate pushdown and why it matters.\\[4pt]
\textbf{A:} Predicate pushdown moves filter conditions as close to the data source as possible.  This reduces the number of rows that flow through expensive operators like joins, dramatically cutting I/O and CPU cost.

\vspace{6pt}
\textbf{Q:} Why is join ordering important?\\[4pt]
\textbf{A:} The intermediate result size depends heavily on join order.  Joining a 1M-row table with a 10-row table first produces at most 10 rows, whereas joining two 1M-row tables first could produce up to $10^{12}$ rows.  The optimiser uses dynamic programming (System R algorithm) to find the cheapest order.

\vspace{6pt}
\textbf{Q:} Compare the Volcano model with the materialisation model.\\[4pt]
\textbf{A:} Volcano (iterator) processes one tuple at a time via \code{next()}, which is memory-efficient but has high per-tuple function-call overhead.  Materialisation produces complete intermediate results, which is simpler but can use excessive memory.  Vectorised execution (DuckDB) is a hybrid: it processes batches of tuples, amortising function-call overhead while staying cache-friendly.
\end{interviewbox}
